\section*{Capítulo \ref{ch:b2:global-change} --- El cambio global y la biosfera}

\begin{solution}{exo:b2:global-change:1}
$\Delta T = \lambda C_{\text{acum}}$ con $\lambda = \qty{1.65}{\celsius}$
por \qty{1000}{GtC}: \qty{1.2}{\celsius}, \qty{1.65}{\celsius},
\qty{3.3}{\celsius}.
\end{solution}

\begin{solution}{exo:b2:global-change:2}
Grados-día: la suma, a lo largo de los días, del exceso de la temperatura
media diaria sobre una base, el \omterm{thm:b2:global-change:phenology}{tiempo térmico} que rige el desarrollo de
las plantas y de los ectotermos. Desajuste: cuando dos especies que
interactúan regulan sus estaciones mediante señales distintas y el
calentamiento desplaza a una y no a la otra --- el papamoscas, guiado por
la duración del día en África, llega en su fecha de siempre y encuentra
que el máximo de orugas, guiado por la temperatura, pasó hace ya dos
semanas.
\end{solution}

\begin{solution}{exo:b2:global-change:3}
Unos \qty{150}{m} cuesta arriba y \qty{150}{km} hacia el polo por grado.
Una especie de la cima no tiene terreno más alto: su isoterma sube por
encima de la montaña y su hábitat desaparece, mientras que el área de cada
banda se reduce con la altitud en la subida.
\end{solution}

\begin{solution}{exo:b2:global-change:4}
El \ensuremath{\mathrm{CO_2}} disuelto forma ácido carbónico, que libera iones
hidrógeno, y el pH baja en proporción al logaritmo de la presión de
\ensuremath{\mathrm{CO_2}}. Los iones hidrógeno adicionales convierten el
carbonato en bicarbonato, de modo que bajan la concentración de carbonato
y el \omterm{thm:b2:global-change:acid}{estado de saturación} del carbonato de calcio. Los organismos que
construyen aragonito --- corales, pterópodos, muchas larvas --- y calcita
--- cocolitóforos, foraminíferos, moluscos --- deben gastar más para
construir sus conchas y, por debajo de la saturación, las pierden.
\end{solution}

\begin{solution}{exo:b2:global-change:5}
\qty{1.5}{\celsius}: $1.5/1.65\times 1000 = \qty{909}{GtC}$ en total,
\qty{209}{GtC} restantes, 21 años. \qty{2}{\celsius}:
\qty{1212}{GtC}, \qty{512}{GtC} restantes, 51 años.
\end{solution}

\begin{solution}{exo:b2:global-change:6}
$\sqrt{2\times 150/0.12} = 50$ días; adelanto de $1.5/0.12 = 12.5$
días.
\end{solution}

\begin{solution}{exo:b2:global-change:7}
\qty{3}{\celsius}: las isotermas suben $3000/6.5 = \qty{460}{m}$; la
banda pasa a \qtyrange{1260}{1860}{m}, pero la montaña termina a
\qty{1800}{m}: se pierden los \qty{60}{m} superiores de la banda y el resto
se sitúa donde el área es $2^{-460/300} = 0.35$ de lo que era --- se han
perdido unos dos tercios del área. \qty{5}{\celsius}: \qty{770}{m}; la
banda estaría en \qtyrange{1570}{2170}{m}; solo existe
\qtyrange{1570}{1800}{m}, un tercio de la anchura de la banda, con
$2^{-770/300} = 0.17$ del área --- menos de una décima parte del área
original: la especie ha desaparecido en la práctica.
\end{solution}

\begin{solution}{exo:b2:global-change:8}
560: pH $7.93$, iones hidrógeno $\times 1.9$; 800: $7.79$, $\times
2.6$; 1000: $7.70$, $\times 3.1$.
\end{solution}

\begin{solution}{exo:b2:global-change:9}
$1 - (1 - h)^{0.25}$: \qty{26}{\%}, \qty{53}{\%}, \qty{82}{\%}. La
pérdida inmediata es menor porque el fragmento restante sigue albergando
poblaciones de la mayoría de las especies; son demasiado pequeñas para
persistir y se extinguen a lo largo de décadas --- la \omterm{thm:b2:global-change:sar}{deuda de extinción}.
\end{solution}

\begin{solution}{exo:b2:global-change:10}
$c = 2\sqrt{0.7\times 20} = \qty{7.5}{km/yr}$; \qty{1000}{km} en unos
130 años. Reducir a la mitad $r$ o $D$ divide $c$ por
$\sqrt 2$, hasta \qty{5.3}{km/yr}; ninguna de las dos lo reduce a la mitad
--- para reducir la velocidad a la mitad hay que dividir por cuatro el
producto $rD$.
\end{solution}

\begin{solution}{exo:b2:global-change:11}
Tasa de fondo: $8\times 10^{6}\times 10^{-6} = 8$ al año, 800 por siglo;
con mil veces más, 8000 al año, \num{800000} por siglo. Diez mil
extinciones documentadas en un siglo son unas diez veces la tasa de fondo,
no mil --- pero la documentación solo abarca el pequeño porcentaje de
especies descritas y seguidas (aves, mamíferos, algunas plantas), entre las
que la tasa es, en efecto, cien veces la de fondo o más; en la mayoría no
descrita, la mayor parte de las extinciones nunca se ven, y la estimación
especies--área es el único asidero. La verdad está entre lo contado y lo
estimado, y por eso el intervalo citado es tan amplio.
\end{solution}

\begin{solution}{exo:b2:global-change:12}
La temperatura es proporcional a las \omterm{prop:b2:global-change:tcre}{emisiones acumuladas}; cuando las
emisiones netas son nulas, el total acumulado deja de crecer y el
calentamiento también. En el \omterm{thm:b2:biogeochemical-cycles:box}{modelo de caja}, con emisiones nulas, el exceso
atmosférico empieza a bajar a medida que actúan los sumideros, lo que
enfriaría --- pero el océano, que sigue calentándose hacia el equilibrio
con el \ensuremath{\mathrm{CO_2}} elevado, sigue calentando la superficie, y los dos
efectos casi se compensan: la temperatura se mantiene donde está durante
siglos. Lo que le falta al argumento: los demás gases de \omterm{prop:b2:global-change:tcre}{efecto
invernadero} (la corta vida del metano hace que su calentamiento baje
deprisa cuando cesan las emisiones), los aerosoles que hoy enmascaran
parte del calentamiento y desaparecen en unas semanas cuando cesan las
emisiones que los producen, y las retroalimentaciones --- el carbono del
permafrost, el decaimiento de los bosques --- que podrían añadir emisiones
propias.
\end{solution}

\begin{solution}{pb:b2:global-change:1}
\textbf{1.} A: $700 + 80\times 10 = \qty{1500}{GtC}$. B: $700 +
\tfrac12\times 50\times 10 = \qty{950}{GtC}$.
\textbf{2.} A: $1.65\times 1.5 = \qty{2.5}{\celsius}$. B:
$1.65\times 0.95 = \qty{1.6}{\celsius}$.
\textbf{3.} 2050. A: \qty{1000}{GtC}, \qty{1.65}{\celsius}. B:
emisiones $10\int_0^{30}(1 - t/50)\,\mathrm{d}t = 10(30 - 9) =
\qty{210}{GtC}$, total de 910, \qty{1.5}{\celsius}.
\textbf{4.} $\lambda\times 100 = \qty{0.165}{\celsius}$ por década,
algo por debajo de los $0.2$ observados (que incluyen los demás gases y
la pérdida del enmascaramiento por los aerosoles).
\textbf{5.} Sí: las \omterm{prop:b2:global-change:tcre}{emisiones acumuladas} dejan de crecer en 2070, y el
calentamiento también, que se mantiene en \qty{1.6}{\celsius} --- la
proporcionalidad significa que la temperatura la fija el total, no el
ritmo.
\textbf{6.} $2/1.65\times 1000 = \qty{1212}{GtC}$; la trayectoria A la
alcanza al cabo de $512/10 = 51$ años, en 2071.
\textbf{7.} $\sqrt{2\times 200/0.1} = 63$ días.
\textbf{8.} Calentamiento respecto de 2020: A, $2.5 - 1.2 =
\qty{1.3}{\celsius}$; B, \qty{0.4}{\celsius}. Adelanto $\Delta T/b$:
A, 13 días; B, 4 días.
\textbf{9.} Intervalo: A, $20 + 13 = 33$ días; B, 24 días.
\textbf{10.} Orugas disponibles del día 15 al día 25 tras la foliación.
B: el ave llega el día 24, dentro de la ventana, por poco. A: el día 33,
ocho días después de la última oruga.
\textbf{11.} El ave se adelanta $3\times 1.3 = 4$ días: intervalo de $33 - 4 =
29$ días --- todavía fuera de la ventana.
\textbf{12.} El ave puede vivir a la nueva temperatura; lo que no puede es
alimentar a sus polluelos cuando el alimento ya ha desaparecido, porque su
señal y la de su presa responden de forma distinta al calentamiento.
\textbf{13.} A: $1.3/6.5\times 1000 = \qty{200}{m}$ hacia arriba y
$1.3/0.65\times 100 = \qty{200}{km}$ hacia el polo. B: \qty{57}{m} y
\qty{57}{km}.
\textbf{14.} En 80 años, el árbol se desplaza \qty{8}{km}: no puede seguir
\qty{200}{km} (A) y tampoco llega a \qty{57}{km} (B) --- los árboles van
con retraso, y sus áreas de distribución estarán fuera de equilibrio con el
clima durante siglos en cualquiera de las dos trayectorias.
\textbf{15.} A: la banda estaría en \qtyrange{1700}{2000}{m}: alcanza
justo la cima y no tiene adónde seguir subiendo. La montaña se estrecha
con la altura, de modo que la misma franja ocupa ya solo $2^{-200/300} =
0.63$ veces su área anterior --- y un calentamiento mayor no deja terreno
alguno al que desplazarla.
\textbf{16.} B: \qty{57}{m} hacia arriba, factor de área $2^{-57/300} = 0.88$: una
pérdida del \qty{12}{\%}.
\textbf{17.} Ocho décadas a \qty{30}{km}: \qty{240}{km}, más que los
\qty{200}{km} necesarios en A: el pez sigue el ritmo (y abandona su antigua
pesquería).
\textbf{18.} En el borde inferior, la especie se encuentra con
competidores, patógenos y depredadores que suben desde abajo y a los que
antes no se enfrentaba, y su propia fisiología, adaptada a la antigua
temperatura, rinde menos que la de ellos: el retroceso es tan competitivo
como fisiológico.
\textbf{19.} A: $\mathrm{pH} = 8.2 - 0.9\log_{10}(600/280) = 7.90$;
iones hidrógeno, $10^{0.30} = 2.0$ veces los de 1750. B: $8.04$; $1.44$
veces.
\textbf{20.} $\Omega = 4/(\text{cociente})$: A, $2.0$; B, $2.8$. Con 2, los
corales siguen calcificando, pero despacio, con esqueletos más débiles y
arrecifes que se erosionan más deprisa de lo que crecen; con $2.8$
construyen, aunque el calor de la trayectoria B sigue blanqueándolos.
\textbf{21.} A: $400\times 0.4^{0.25} = 318$, 82 perdidas. B:
$400\times 0.8^{0.25} = 378$, 22 perdidas.
\textbf{22.} A: $318\times(1 - 0.05\times 1.3) = 298$. B:
$378\times(1 - 0.05\times 0.4) = 371$.
\textbf{23.} A: $1 - 298/400 = \qty{26}{\%}$. B: \qty{7}{\%}.
\textbf{24.} En ambas, la tala: 82 frente a 20 especies en A, 22 frente a 7
en B --- el término de calentamiento del modelo es modesto; en realidad,
ambos interactúan, puesto que un bosque fragmentado no puede desplazar su
área de distribución.
\textbf{25.} Trayectoria A: \qty{2.5}{\celsius}, pH de $7.90$, una cuarta
parte de las especies de árboles perdidas. Trayectoria B:
\qty{1.6}{\celsius}, pH de $8.04$, \qty{7}{\%} perdidas.
\end{solution}
